\scanpage{5}%
\mspage{\handwritten{8}}%
phisme $\phi^{-1}(a)\handwritten{\simeq}M_a$, transformant $(e,a)$ en $e$, \margin{\handwritten{e}}c'est une égalité si et seulement si $M_a$ est \inserted{\handwritten{régulier en \deleted{\unsure{$a$}}$e$ i.e.}} lisse sur $k$ en $e$ i.e. $\phi^{-1}(a)$ est lisse sur $k$ en $(e,a)$. Si \deleted{\unsure{donc}} $\handwritten{N}=\handwritten{M}_a$, alors les membres extrêmes de la suite d'inégalités précédente sont égaux, donc tous les termes sont égaux, ce qui implique que $\phi$ est universellement ouvert en $(e,a)$ et que sa fibre en ce point est lisse sur $k$. Comme $V$ est \deleted{réduit \unsure{en}} normal donc réduit en $e$, il \deleted{résulte \unsure{d'un}}\inserted{s'ensuit, (\handwritten{p}ar un} théorème de EGA IV) que $\phi$ est lisse \deleted{\unsure{sur}} en $(e,a)$, ce qui prouve 1.4. La démonstration invoque certains résultats assez délicats de \inserted{la} théorie des schémas, \deleted{\unsure{mais}} mais nous n'au\handwritten{\unsure{t|r}}ons pas besoin de 1.4. pour la suite.

Les conditions de 1.3. étaient relatives au choix d'un point $a$ de $\mathbb{W}(k)$. On conclut de 1.3. l'énoncé suivant, indépendant d'un tel choix :

\uline{Corollaire} 1.5. Soient $G,V,W$ sur $k$ comme dans 1.3. \deleted{Considérons les} et supposons $V$ génériquement \deleted{réduit. Alors les seul}\inserted{\handwritten{séparable sur $k$.}} Considérons les conditions suivantes:

(i) \deleted{\unsure{$\phi$}\illegible{1}} $\handwritten{\psi}:X\handwritten{\to}V$ est génériquement non ram\deleted{\unsure{ff|if}}\handwritten{if}ié.

(i bis) $\handwritten{\psi}:X\handwritten{\to}V$ est génériquement étale.

(i\deleted{\unsure{i}} \deleted{bis}\handwritten{ter}) $\handwritten{\psi}:X\handwritten{\to}V$ est génériquement lisse, $N$ est lisse sur $k$ et l'inégalité (2) une égalité .

(ii) \deleted{Il existe $a$ \unsure{$\in$} $W(k)$ tel que} $N$ est lisse sur $k$, et il existe $a\handwritten{\in}W(k)$ tel que $\handwritten{M}=\handwritten{N}_a$.

(iii) $\handwritten{\psi}:X\handwritten{\to}V$ est génériquement quasi-fini, et il existe $a\handwritten{\in}W(k)$ tel que $\handwritten{M}=\handwritten{N}_a$.

(iv) Il existe $a\handwritten{\in}W(k)$ tel que $M_a$ et $N$ coïncident au voisinage de $e$\handwritten{.}

Alors on a les implications
\[
\text{(ii)}\handwritten{\Leftrightarrow}\text{(iii)}\handwritten{\Leftrightarrow}\text{(iv)}
\handwritten{\Longrightarrow}\text{(i)}\handwritten{\Leftrightarrow}\text{(i bis)}\handwritten{\Leftrightarrow}\text{(i ter)}\ ,
\]
et si $k$ est algébriquement clos, toutes les conditions envisagées sont
% archive sheet number: 4
